\ssscp{Bat/flying fox}{C-10-01-02}
Words meaning ‘flying fox’
and ‘bat’ in \tsc{Greater West Bomberai} languages are reconstructed ultimately as *mandəl
\citep[493]{UsherSchapper2022},
 which is remarkably close in form to **mantəR/**mandəR ‘cuscus’.
Both bats and cuscus are edible mammals dwelling in the treetops.
We also note that some languages of eastern Seram have an explicit
historical link between their words for ‘cuscus’
and flying creatures, including bats.
Consider \ili{Masiwang} \it{mani} ‘cuscus’ (<*manuk ‘bird,
flying creature’; ultimate *u > \it{i} is regular,
as shown in Tables \ref{tab:11-11}
and \ref{tab:11-12}); and \ili{Bati} \it{manuk faniki} ‘cuscus’ (<*manuk ‘bird,
flying creature’; *paniki ‘fruit bat, flying fox’).\footnote{
		\citet[269]{Schapper2011} provided a Proto-TAP reconstruction *mader ‘bat’.
		The final consonant of the reconstruction is based on \ili{Sawila} \it{madiirang} ‘bat’.
		But that form is incorrect.
		\ili{Sawila} \it{madiirang} means ‘sharp’.
		One item up in the word list is \it{madiiku} ‘bat,
		(Indonesian \it{kelelawar})’.}
At this point we are merely observing the remarkable similarities
in form, and possible links in meaning,
which are too close to simply dismiss out of hand.

\newpage
\noindent{\wg\st{2pt}\begin{xltabular}{\linewidth}
{lllll>{\raggedright\arraybackslash}X}
\lsptoprule
%Language	&	ISO	&	Form	&	gloss	&	Etymon	&	Source	\\	\midrule
%\endfirsthead
Language	&	ISO	&	Form	&	gloss	&	Etymon	&	Source	\\	\midrule	\hhline{~}
\endhead
\ili{Iha}	&	ihp	&	manda;	&	bat	&	*mandəl;	&	SV \citeyear[122]{SmitsVoorhoeve1998}	\\	\hhline{~}
	&		&	mbayɛr	&		&	\it{misc.}	&	\cite{Donohue2010}	\\
\ili{Mbaham}	&	bdv	&	mbayé	&	bat	&	\it{misc.}	&	\cite[63]{Cottet2015}	\\
\ili{Kalamang}	&	kgv	&	kuek	&	fruit bat	&	\it{misc.}	&	\cite[467]{Visser2022}	\\
\ili{Makasae}	&	mkz	&	siri	&	bat	&	\it{misc.}	&	\cite{Klamer2002}	\\
\ili{Makalero}	&	mjb	&	hiriʔ	&	bat	&	\it{misc.}	&	\cite[556]{Huber2011}	\\
\ili{Bunak}	&	bfn	&	zupal	&	bat	&	\it{misc.}	&	\cite{Klamer2002}	\\
\ili{Fataluku}	&	ddg	&	mat͡sa	&	bat	&	*mandəl	&	\cite[189]{Heston2015}	\\
\ili{Oirata}	&	oia	&	maʈa-mailur	&	large bat	&	*mandəl	&	\cite{DeJosselinDeJong1937} cited in \cite[493]{UsherSchapper2022}	\\
\ili{Wersing}	&	kvw	&	məduku	&	bat	&	*mandəl	&	Johnny Banamtuan pers. comm. Oct.2020	\\
\ili{Sawila}	&	swt	&	madiiku 	&	bat	&	*mandəl	&	Johnny Banamtuan pers. comm. Oct.2020	\\
\ili{Kamang}	&	woi	&	maˈtei	&	bat	&	*mandəl	&	\cite[87]{SchapperManimau2011} (Dialect variants: \it{mate, matii})\\
\ili{Abui}	&	abz	&	maˈrɛl	&	bat	&	*mandəl	&	\cite[85]{KratochvilDelpada2008}	\\
\ili{Kafoa}	&	kpu	&	marel	&	bat	&	*mandəl	& \cite[493]{UsherSchapper2022}	\\
\ili{Klon}	&	kyo	&	mədel	&	bat	&	*mandəl	&	\cite[269]{Schapper2011}	\\
\ili{Teiwa}	&	twe	&	medi	&	bat	&	*mandəl	&	\cite[35]{KlamerSir2011}	\\
\ili{Klamu}\footnotemark 	&	nec	&	[ˈmărˑɑ]	&	bat	&	*mandəl	&	Grimes fieldnotes (see \srf{02-04-06})	\\
\ili{Papuna}	&	{-}{-}{-}	&	marie	&	bat	&	*mandəl	&	\cite{Delpada2011}	\\
\ili{Mauta}	&	lev	&	madːe	&	bat	&	*mandəl	&	\cite[92]{HoltonKoly2008} (k.o. slightly furry bat)	\\
\ili{Tubbe}	&	lev	&	madːe	&	bat	&	*mandəl	&	\cite[92]{HoltonKoly2008} (k.o. slightly furry bat)	\\
\ili{Lamma}	&	lve	&	madːal	&	flying fox	&	*mandəl	&	\cite[92]{HoltonKoly2008} (Indonesian \it{kalung})	\\
\lspbottomrule
\end{xltabular}}
\footnotetext{
	\citet[493]{UsherSchapper2022} list Nedebang \it{marːa} ‘bat’.
	Nedebang is an alternate name in the literature for \ili{Klamu}.}

%\noindent{\wg\st{2pt}\begin{xltabular}{\linewidth}
%{lllll>{\raggedright\arraybackslash}X}
%\lsptoprule
%
%\lspbottomrule
%\end{xltabular}}